\section{Cuarta fase: escenarios, incertidumbre y elección de política}\label{sec:escenarios}

\begin{idea}{Dos preguntas que requieren respuestas distintas}
\textbf{Escenario:} ¿qué cambia si la economía, la normativa o el suministro son
distintos? \textbf{Incertidumbre dentro del escenario:} aun bajo esos supuestos,
¿cuánto pueden variar los meses? Las sensibilidades examinan además si una
hipótesis poco identificada cambia la decisión. No son tres nombres para la misma cosa.
\end{idea}

\subsection{De una historia a parámetros concretos}
El central mantiene las condiciones económicas del cierre de A012. Los extremos
utilizan un modelo macroeconómico auxiliar para obtener un único cociente entre
mercado bajo el escenario y mercado auxiliar central. Ese cociente se aplica a
la previsión elegida: no se sustituye el central por otro modelo ni se suma una
penalización económica genérica adicional.

\begin{table}[htbp]\centering
\caption{Principales supuestos A013. Las variaciones de indicadores se miden
respecto al cierre de A012; no son variaciones de ventas.}\label{tab:escenarios-param}
\small
\begin{tabularx}{\linewidth}{@{}p{3.6cm}YYY@{}}\toprule
Factor & Pesimista & Central & Optimista \\\midrule
Paro & +2 puntos porcentuales & Sin cambio & $-1$ punto porcentual \\
Confianza & $-16{,}26$ puntos de índice & Sin cambio & +19,14 puntos de índice \\
Tipo de crédito & +0,546 puntos porcentuales & Sin cambio & $-0{,}546$ puntos porcentuales \\
Devolución del adelanto & 6 meses & 3 meses & 1 mes \\
Ayuda pública & No & No & Abril--diciembre \\
Producto & Central neutro; retraso en sensibilidad & Efecto adicional cero & +16,8\,\% hipotético durante 6 meses, con normalización previa \\
Turismo & Cambio bajo de viajeros & Sin cambio & Cambio alto de viajeros \\
Combustible residual & $-3\,\%$ de demanda & Cero & $+3\,\%$ de demanda \\
CO$_2$ & Prima del catálogo por venta & Sin prima & Sin prima \\
Nueva importación & +500 \euro/coche por arancel & Sin recargo & $-250$ \euro/coche por divisa \\
Suministro MSUV & Pérdida de producción en marzo; llegada junio & Sin interrupción & Sin interrupción \\\bottomrule
\end{tabularx}
\fuente{\archivo{resultados/escenarios/tablas/parametros_escenarios_a013.csv}}
\end{table}

Los choques de confianza y viajeros se apoyan en percentiles de cambios históricos
a nueve meses; el del crédito, en la mediana absoluta de esos cambios. El paro
sigue la propuesta del enunciado. Turismo afecta al canal de alquiler durante
marzo--junio, ponderado por su peso en los clientes de Kaido, no por el peso del
alquiler en todo el mercado.

El combustible, el arancel y el cambio no tienen una base suficiente para estimar
sus tamaños causalmente. Sus importes son \textbf{pruebas de estrés explícitas},
contrastadas con valores cero y dobles. El coste arancelario no se calcula como
un porcentaje arbitrario del PVP. La divisa sólo cambia el coste de las nuevas
importaciones; no revaloriza gratuitamente los coches iniciales.

\subsection{En qué orden se aplican los efectos}
La composición importa. Primero se aplican efectos sobre mercado y posición,
turismo y combustible. Después se descuenta el adelanto expresado en coches:
\begin{equation}\label{eq:demanda-escenario}
 \widehat D^{(e)}_{s,t}
 =\max\!\left(0,\ \widehat D^{(e),\mathrm{previa}}_{s,t}(1+c_e)
                       -\frac{A_s}{m_e}\,\mathbf1_{t\leq m_e}\right).
\end{equation}
$c_e$ es el efecto residual del combustible; $A_s$, el adelanto A012; y $m_e$,
el número de meses de devolución. La misma cantidad central $A_s$ se solicita
en los tres escenarios. Multiplicar después por $1+c_e$ alteraría también el
adelanto absoluto, que no debe convertirse en una segunda elasticidad.

En el optimista, concentrar la devolución en enero supera la demanda prevista
de algunos segmentos. El límite a cero deja \num{235.3} coches del adelanto sin
descontar en la integración final. Se declara esa parte no aplicada; no se crea
demanda negativa ni se traslada silenciosamente a febrero. Es una limitación de
la combinación de hipótesis, no un error de redondeo.

\subsection{Generar años posibles sin perder la dependencia entre meses}
Una réplica produce 18 meses de demanda realizada: A013 y seis meses de
continuidad. La realización se genera mediante
\begin{equation}\label{eq:ruido}
 D^{(r,e)}_{s,t}=\operatorname{redondear}
 \left[\widehat D^{(e)}_{s,t}
       \exp\!\left(\sigma_{\log,s}Z^{(r)}_{s,t}
                         -\frac{\sigma_{\log,s}^{2}}{2}\right)\right].
\end{equation}
El factor exponencial es positivo. Si $Z$ es normal de media cero y varianza uno,
restar $\sigma_{\log,s}^{2}/2$ hace que el factor tenga media uno antes del
redondeo. No significa que cada réplica coincida con la previsión ni que su
mediana sea igual a su media.

El ruido contiene una parte común y una específica:
\begin{align}
 Z_{s,t}&=\sqrt\rho\,U_t+\sqrt{1-\rho}\,W_{s,t},\\
 U_t&=\phi U_{t-1}+\sqrt{1-\phi^2}\,\eta_t,
\end{align}
con una ecuación análoga para $W$. Las innovaciones $\eta$ son normales
independientes; $\phi$ introduce persistencia temporal y $\rho$ co-movimiento
entre segmentos. En la ejecución, $\phi\simeq0{,}210$ y $\rho\simeq0{,}190$.
No se asume que un mes alto se compense necesariamente con uno bajo.

\tabla{incertidumbre.tex}{Dispersión de errores por segmento. Son parámetros
adimensionales; $\sigma_{3,\mathrm{rel}}$ corresponde a errores acumulados en tres meses.}
{tab:incertidumbre}{\archivo{resultados/escenarios/tablas/parametros_incertidumbre_a013.csv}}

Los parámetros proceden de los errores de selección en segmentos maduros,
separando años al calcular ventanas y correlaciones. El logaritmo de error usa
un suavizado de 0,5 coche para tratar ceros. CSUV recibe el proxy ampliado. El
valor muy elevado de la dispersión relativa de S revela una serie pequeña e
inestable: no debe interpretarse como una precisión demostrada para ese nicho.

\subsection{Qué ocurre dentro de cada réplica}
La secuencia se repite sin cambiar las reglas de contabilidad:
\begin{enumerate}
  \item Se parte del mismo stock y del mismo plan inicial para los tres escenarios.
  \item Llegan los compromisos del mes, respetando las interrupciones conocidas.
  \item Se atiende la demanda con FIFO y, sólo en la alternativa correspondiente,
  con transfers disponibles.
  \item Se anotan ventas, faltantes, edades, stock y costes.
  \item Se observan ventas y disponibilidad, se revisa la previsión y se decide
  únicamente la llegada de dentro de tres meses.
  \item Se continúa hasta junio de A014 para examinar los efectos de cierre.
\end{enumerate}

La falta de una pieza en fábrica en marzo afecta a la llegada de junio de MSUV.
No borra coches que ya estaban viajando. La producción perdida no se recupera
automáticamente después; hacerlo requeriría otra hipótesis de capacidad.

Todos los escenarios utilizan las mismas innovaciones aleatorias al compararse.
Esto no impone probabilidades a los escenarios: reduce ruido al comparar políticas.
La semilla 42 permite reproducir las trayectorias.

\subsection{Cómo leer media, mediana y P10--P90}
\begin{description}[style=nextline]
\item[Media.] Promedio de las réplicas; es la magnitud usada para la contribución
esperada bajo un escenario.
\item[Mediana.] Mitad de las réplicas por encima y mitad por debajo.
\item[P10 y P90.] Aproximadamente el 10\,\% de las réplicas queda por debajo del
primero y el 10\,\% por encima del segundo. El intervalo intermedio contiene
aproximadamente el 80\,\% de las realizaciones de la simulación.
\end{description}
No son un intervalo de confianza para una verdad desconocida ni una probabilidad
de que se produzca el escenario. Para el total se suman los segmentos
\emph{dentro de cada réplica} y después se calculan sus percentiles. La suma de
medianas de segmentos no es, en general, la mediana del total.

\subsection{Elegir una política: margen sujeto a servicio}
El servicio de un segmento en una réplica es la proporción de su demanda anual
que se ha podido atender:
\begin{equation}\label{eq:servicio}
 F_s^{(r)}=\frac{\sum_{t\in A013}V_{s,t}^{(r)}}
                 {\sum_{t\in A013}D_{s,t}^{(r)}}.
\end{equation}
La restricción acordada es que la \textbf{media de $F_s$ en el central} alcance
el 95\,\% en los segmentos marginalmente viables. No es un 95\,\% de
probabilidad de pasar el año sin roturas ni un requisito en cada escenario.

Se prueban seis niveles de protección:
$z\in\{0;\,0{,}84;\,1{,}28;\,1{,}65;\,2{,}33;\,3\}$.
El piloto usa 200 réplicas. Como no se impone una capacidad compartida de fábrica
no cuantificada en el caso, cada segmento puede escoger su nivel.

Entre políticas factibles se busca mayor contribución media central. Cuando una
opción con menor colchón queda a no más de dos errores estándar de Monte Carlo
de la mejor, se prefiere la menor protección. El error estándar se calcula sobre
las diferencias emparejadas:
\begin{equation}
 \mathrm{EE}(\overline{\Delta R})
 =\frac{\operatorname{sd}(R^{(r)}_{\mathrm{mejor}}-R^{(r)}_z)}{\sqrt N}.
\end{equation}
Esta tolerancia depende de la precisión computacional de la comparación, no
de una cantidad fija idéntica para todos los segmentos. Tampoco mide toda la
incertidumbre económica del modelo.

La política elegida se comprueba con 1.000 réplicas. El código permite aumentar
el colchón si esa verificación descubre un incumplimiento; en la ejecución
presentada no fue necesario cambiar los valores elegidos por el piloto.

\tabla{politica.tex}{Política seleccionada y servicio anual central. La
probabilidad indicada es la frecuencia simulada de $F_s<95\,\%$, condicionada al central.}
{tab:politica}{\archivo{resultados/escenarios/tablas/politica_recomendada_a013.csv}}

SSUV alcanza un servicio medio del 95,4\,\%, pero en el 39,1\,\% de las
réplicas no llega al 95\,\%. No hay contradicción: la restricción está formulada
sobre la media. Su P10 del 89,3\,\% expresa un riesgo que la media oculta.

\begin{figure}[htbp]\centering
\begin{tikzpicture}
\begin{axis}[width=.94\linewidth,height=6.7cm,xmin=-.4,xmax=6.4,ymin=80,ymax=101,
  xtick={0,1,2,3,4,5,6},xticklabels={B,C,MD,SSUV,CSUV,MSUV,S},
  ylabel={Demanda anual atendida (\%)},xlabel={Segmento},
  legend style={at={(.5,-.24)},anchor=north,legend columns=3}]
\addplot[azul,only marks,mark=*,mark size=2.4pt]
  table[col sep=comma,x expr=\coordindex,y expr=100*\thisrow{servicio_central}]
  {datos/politica_recomendada_a013.csv};\addlegendentry{Media central}
\addplot[gris,only marks,mark=square*,mark size=2.1pt]
  table[col sep=comma,x expr=\coordindex,y expr=100*\thisrow{p10_servicio}]
  {datos/politica_recomendada_a013.csv};\addlegendentry{P10 central}
\addplot[azul,dashed,domain=-.4:6.4]{95};\addlegendentry{Objetivo en media}
\end{axis}\end{tikzpicture}
\caption{Cumplir el objetivo medio no elimina años desfavorables.
El eje vertical está ampliado al intervalo 80--101\,\%; son puntos, no barras
de cantidades con origen recortado.}
\label{fig:servicio}\fuente{política recomendada, 1.000 réplicas centrales}
\end{figure}

\begin{idea}{Alcance de la elección}
Se recomienda la mejor opción entre las políticas comparadas bajo el objetivo
declarado. No se demuestra un óptimo global ni una política óptima para cualquier
probabilidad de escenarios. La malla, la ausencia de capacidad compartida y la
forma de los errores son parte del experimento.
\end{idea}
\FloatBarrier
